\documentclass[11pt,a4paper]{article}
\usepackage{fontspec,amsmath,amssymb,amsthm,microtype}
\setmainfont{lmroman10-regular.otf}[BoldFont=lmroman10-bold.otf,ItalicFont=lmroman10-italic.otf,BoldItalicFont=lmroman10-bolditalic.otf]
\usepackage{xcolor}
\usepackage[most]{tcolorbox}
\usepackage{mathrsfs}
\providecommand{\lt}{<}
\providecommand{\gt}{>}
\setmonofont{DejaVuSansMono.ttf}[BoldFont=DejaVuSansMono-Bold.ttf,ItalicFont=DejaVuSansMono-Oblique.ttf,BoldItalicFont=DejaVuSansMono-BoldOblique.ttf,Scale=0.86]
\newfontfamily\leanfallback{STIXTwoMath-Regular.otf}[Scale=0.86]
\definecolor{leanrule}{gray}{0.62}
\definecolor{leanback}{gray}{0.962}
\definecolor{leantitle}{gray}{0.22}
\definecolor{leanmuted}{gray}{0.38}
\newtcolorbox{leanbox}[1]{enhanced,breakable,colback=leanback,colframe=leanrule,
  boxrule=0.4pt,arc=0.6mm,left=2.4mm,right=2.4mm,top=1.8mm,bottom=1.8mm,
  before skip=0.8em,after skip=1em,fonttitle=\small\bfseries,coltitle=leantitle,
  colbacktitle=leanback,titlerule=0.3pt,title={#1}}
\newlength{\leancharwidth}
\newenvironment{leancode}{\par\smallskip\begingroup\ttfamily\footnotesize\raggedright
  \settowidth{\leancharwidth}{M}\setlength{\parskip}{0pt}\setlength{\parindent}{0pt}
  \frenchspacing}{\par\endgroup\smallskip}
\newcommand{\leanline}[2]{\hangindent=\dimexpr#1\leancharwidth+4\leancharwidth\relax
  \hangafter=1\noindent\hspace*{#1\leancharwidth}#2\par}
\newcommand{\leansource}[1]{\par{\footnotesize\color{leanmuted}#1}\par}
\newcommand{\leanlabel}[1]{\par\smallskip\noindent{\small\bfseries #1}\par}
\usepackage[margin=25mm]{geometry}
\usepackage[unicode,hidelinks]{hyperref}
\hypersetup{pdfauthor={LangevinLean Project}}
\setlength{\emergencystretch}{3em}
\setlength{\parskip}{0.4em}
\setlength{\parindent}{0pt}
\allowdisplaybreaks[2]
\raggedbottom
\title{\texorpdfstring{Curvature, step size and the potential contract}{Curvature, step size and the potential contract}}
\author{LangevinLean Project}
\date{}
\begin{document}\maketitle
Draft mathematical exposition; these notes may contain mistakes.


Let \(0<\alpha\le\beta\), let \(d\) be a nonnegative integer and let \(h>0\) satisfy \(h\beta\le1\). Define \(\kappa=\beta/\alpha\), \(q=1-\alpha h\) and \(B=6h\sqrt{\kappa\beta d}\). These are numerical parameters; no invariant law or spectral representation is included in their definition.



Earlier results used here: \cite[Section~3.10]{LangevinLean2026ULASC}.

\section{\texorpdfstring{The numerical parameters and their elementary bounds}{The numerical parameters and their elementary bounds}}


\textbf{Claim.} Under the displayed restrictions, \(\beta>0\), \(\kappa\ge1\), \(0\le q<1\), \(h\le1/\beta\), and \(B\ge0\).



\textbf{Proof.} Since \(\beta\ge\alpha>0\), division by \(\alpha\) preserves \(\beta\ge\alpha\) and gives \(\kappa\ge1\). Multiplication by \(h\ge0\) gives \(\alpha h\le\beta h\le1\), so \(q\ge0\). The strict inequality \(\alpha h>0\) gives \(q<1\). Dividing \(h\beta\le1\) by \(\beta>0\) gives the step restriction. Finally \(h>0\) and the nonnegativity of the square root imply \(B\ge0\).

\begin{leanbox}{Lean: \texttt{Parameters}, \texttt{Parameters.\allowbreak{}condition}, \texttt{Parameters.\allowbreak{}contraction} and 7 more}
\leansource{\texttt{Parameters} \quad WassersteinMixing/\allowbreak{}Parameters.lean, lines 19–27}
\begin{leancode}
\leanline{0}{structure\ Parameters\ where}
\leanline{2}{alpha\ :\ ℝ}
\leanline{2}{beta\ :\ ℝ}
\leanline{2}{step\ :\ ℝ}
\leanline{2}{dimension\ :\ ℕ}
\leanline{2}{alpha\_pos\ :\ 0\ <\ alpha}
\leanline{2}{alpha\_le\_beta\ :\ alpha\ ≤\ beta}
\leanline{2}{step\_pos\ :\ 0\ <\ step}
\leanline{2}{step\_beta\_le\ :\ step\ *\ beta\ ≤\ 1}
\end{leancode}
\leansource{\texttt{Parameters.\allowbreak{}condition} \quad WassersteinMixing/\allowbreak{}Parameters.lean, lines 30–30}
\begin{leancode}
\leanline{0}{def\ Parameters.condition\ (P\ :\ Parameters)\ :\ ℝ\ :=\ P.beta\ /\ P.alpha}
\end{leancode}
\leansource{\texttt{Parameters.\allowbreak{}contraction} \quad WassersteinMixing/\allowbreak{}Parameters.lean, lines 33–33}
\begin{leancode}
\leanline{0}{def\ Parameters.contraction\ (P\ :\ Parameters)\ :\ ℝ\ :=\ 1\ -\ P.alpha\ *\ P.step}
\end{leancode}
\leansource{\texttt{Parameters.\allowbreak{}bias} \quad WassersteinMixing/\allowbreak{}Parameters.lean, lines 36–37}
\begin{leancode}
\leanline{0}{def\ Parameters.bias\ (P\ :\ Parameters)\ :\ ℝ\ :=}
\leanline{2}{6\ *\ P.step\ *\ Real.sqrt\ (P.condition\ *\ P.beta\ *\ P.dimension)}
\end{leancode}
\leansource{\texttt{Parameters.\allowbreak{}beta\_\allowbreak{}pos} \quad WassersteinMixing/\allowbreak{}Parameters.lean, lines 39–40}
\begin{leancode}
\leanline{0}{theorem\ Parameters.beta\_pos\ (P\ :\ Parameters)\ :\ 0\ <\ P.beta}
\end{leancode}
\leansource{\texttt{Parameters.\allowbreak{}condition\_\allowbreak{}ge\_\allowbreak{}one} \quad WassersteinMixing/\allowbreak{}Parameters.lean, lines 42–43}
\begin{leancode}
\leanline{0}{theorem\ Parameters.condition\_ge\_one\ (P\ :\ Parameters)\ :\ 1\ ≤\ P.condition}
\end{leancode}
\leansource{\texttt{Parameters.\allowbreak{}contraction\_\allowbreak{}nonneg} \quad WassersteinMixing/\allowbreak{}Parameters.lean, lines 45–48}
\begin{leancode}
\leanline{0}{theorem\ Parameters.contraction\_nonneg\ (P\ :\ Parameters)\ :\ 0\ ≤\ P.contraction}
\end{leancode}
\leansource{\texttt{Parameters.\allowbreak{}contraction\_\allowbreak{}lt\_\allowbreak{}one} \quad WassersteinMixing/\allowbreak{}Parameters.lean, lines 50–52}
\begin{leancode}
\leanline{0}{theorem\ Parameters.contraction\_lt\_one\ (P\ :\ Parameters)\ :\ P.contraction\ <\ 1}
\end{leancode}
\leansource{\texttt{Parameters.\allowbreak{}step\_\allowbreak{}le\_\allowbreak{}inv\_\allowbreak{}beta} \quad WassersteinMixing/\allowbreak{}Parameters.lean, lines 54–55}
\begin{leancode}
\leanline{0}{theorem\ Parameters.step\_le\_inv\_beta\ (P\ :\ Parameters)\ :\ P.step\ ≤\ 1\ /\ P.beta}
\end{leancode}
\leansource{\texttt{Parameters.\allowbreak{}bias\_\allowbreak{}nonneg} \quad WassersteinMixing/\allowbreak{}Parameters.lean, lines 57–59}
\begin{leancode}
\leanline{0}{theorem\ Parameters.bias\_nonneg\ (P\ :\ Parameters)\ :\ 0\ ≤\ P.bias}
\end{leancode}
\leanlabel{How the Lean proof follows the argument}
\texttt{Parameters} records all numerical hypotheses. \texttt{condition}, \texttt{contraction} and \texttt{bias} define their literal expressions. The remaining declarations divide only by a proved positive curvature constant; the contraction inequalities use the supplied step restriction and positive step.
\end{leanbox}

\section{\texorpdfstring{The first-order potential contract and the original Hessian bounds}{The first-order potential contract and the original Hessian bounds}}


\textbf{Definition.} On a real inner product space, the potential contract requires \(V\in C^2\), \(\alpha\)-strong convexity and a \(\beta\)-Lipschitz gradient:

\[V(y)\ge V(x)+\langle\nabla V(x),y-x\rangle+\frac\alpha2\|y-x\|^2,\qquad\|\nabla V(y)-\nabla V(x)\|\le\beta\|y-x\|.\]

The original source Hessian bounds imply this contract with exactly the same constants, as proved in the Hessian-assumption argument. The contract contains no invariant law, semigroup, spectral basis or comparison conclusion; those objects and properties are constructed separately.

\begin{leanbox}{Lean: \texttt{PotentialAssumptions}}
\leansource{\texttt{PotentialAssumptions} \quad WassersteinMixing/\allowbreak{}Parameters.lean, lines 63–67}
\begin{leancode}
\leanline{0}{structure\ PotentialAssumptions\ \{H\ :\ Type*\}\ [NormedAddCommGroup\ H]}
\leanline{4}{[InnerProductSpace\ ℝ\ H]\ [CompleteSpace\ H]\ (P\ :\ Parameters)\ (V\ :\ H\ →\ ℝ)\ :\ Prop\ where}
\leanline{2}{regularity\ :\ ContDiff\ ℝ\ 2\ V}
\leanline{2}{strongConvexity\ :\ StrongConvexOn\ Set.univ\ P.alpha\ V}
\leanline{2}{gradientLipschitz\ :\ LipschitzWith\ ⟨P.beta,\ P.beta\_pos.le⟩\ (gradient\ V)}
\end{leancode}
\leanlabel{How the Lean proof follows the argument}
\texttt{Potential\allowbreak{}Assumptions} is a proposition-valued structure containing \texttt{Cont\allowbreak{}Diff \ensuremath{\mathbb R} 2 V}, \texttt{Strong\allowbreak{}Convex\allowbreak{}On} and a \texttt{Lipschitz\allowbreak{}With} gradient condition. None of its fields asserts an invariant measure, a generator realization or a spectral basis.
\end{leanbox}

\begin{thebibliography}{ULA-LSI}
\bibitem[ULA-SC]{LangevinLean2026ULASC}
LangevinLean Project. \emph{Unadjusted Langevin under Strong Convexity}. LL-01, 2026. Draft report, content revision d34ead06987c. \url{https://odakyildiz.github.io/LangevinLean/docs/pdf/UnadjustedLangevin/FullProof.pdf}.
\end{thebibliography}

\end{document}
